\begin{table}[p]\centering
\caption{Changes in the equivalent savings rate from additional leverage}\label{tab:exposure-value}
\begin{minipage}{\linewidth}\small The table reports the change in the equivalent savings rate, in percentage points, from raising the gross exposure of each strategy. A negative change means that the household needs to save less. Each comparison is based on 10,000 bootstrap simulations. The last column draws returns from 1970 to 2025, after the 1968 separation of the free London gold market from the official price.\end{minipage}
\medskip

{\setlength{\tabcolsep}{4pt}\small
\begin{tabular}{lrr}\toprule
Gross exposure & \shortstack{Full sample\\ (1927--2025)} & \shortstack{Post-1970 sample\\ (1970--2025)} \\ \midrule
\multicolumn{3}{c}{Panel A: All-equity strategy} \\
100$\to$125 & $-$1.4 & $-$0.9 \\
100$\to$175 & $-$0.4 & +2.5 \\
175$\to$200 & +6.0 & +17.3 \\
\midrule
\multicolumn{3}{c}{Panel B: Proportional strategy} \\
100$\to$125 & $-$2.9 & $-$2.3 \\
100$\to$175 & $-$6.5 & $-$5.3 \\
175$\to$200 & $-$1.1 & $-$0.9 \\
\midrule
\multicolumn{3}{c}{Panel C: Risk parity strategy} \\
100$\to$125 & $-$3.5 & $-$2.6 \\
100$\to$175 & $-$8.0 & $-$6.0 \\
175$\to$200 & $-$1.5 & $-$1.2 \\
\bottomrule\end{tabular}}
\end{table}
